\documentclass{beamer}
\usepackage[utf8]{inputenc}
\usepackage[T1]{fontenc}
\usepackage{graphicx}
\usepackage{amsmath}
\usepackage{amssymb}
\usepackage{listings}
\usepackage{xcolor}
\usepackage{booktabs}

% Theme configuration
\usetheme{Madrid}
\usecolortheme{default}

% Color configuration
\definecolor{codegreen}{rgb}{0,0.6,0}
\definecolor{codegray}{rgb}{0.5,0.5,0.5}
\definecolor{codepurple}{rgb}{0.58,0,0.82}
\definecolor{backcolour}{rgb}{0.95,0.95,0.92}

% Python code style
\lstdefinestyle{pythonstyle}{
    backgroundcolor=\color{backcolour},   
    commentstyle=\color{codegreen},
    keywordstyle=\color{blue},
    numberstyle=\tiny\color{codegray},
    stringstyle=\color{codepurple},
    basicstyle=\ttfamily\footnotesize,
    breakatwhitespace=false,         
    breaklines=true,                 
    captionpos=b,                    
    keepspaces=true,                 
    numbers=left,                    
    numbersep=5pt,                  
    showspaces=false,                
    showstringspaces=false,
    showtabs=false,                  
    tabsize=2,
    language=Python
}

\lstset{style=pythonstyle}

% Document information
\title{Clustering}
\subtitle{Unsupervised Learning}
\author{Meyssem Bouzaien}
\institute{Academic Year: 2025 - 2026}
\date{}

\begin{document}

% Title page
\frame{\titlepage}

% Outline
\begin{frame}
\frametitle{Course Outline}
\tableofcontents
\end{frame}

% ==========================================
% SECTION 1: INTRODUCTION
% ==========================================
\section{What is Clustering?}

\begin{frame}
\frametitle{Reminder: Supervised vs Unsupervised}
\begin{center}
\includegraphics[width=\textwidth]{images/01_supervise_vs_non_supervise.png}
\end{center}
\end{frame}

\begin{frame}
\frametitle{Supervised vs Unsupervised Learning}
\begin{columns}
\begin{column}{0.48\textwidth}
\textbf{Supervised}
\begin{itemize}
    \item Labels provided
    \item We know what we're looking for
    \item Classification, Regression
    \item Example: Spam/Not spam
\end{itemize}
\end{column}

\begin{column}{0.48\textwidth}
\textbf{Unsupervised}
\begin{itemize}
    \item No labels
    \item Discovering structures
    \item Clustering, Dim. reduction
    \item Example: Customer groups
\end{itemize}
\end{column}
\end{columns}

\vspace{1em}
\begin{alertblock}{Key Difference}
Unsupervised = No "right answer" to learn from!
\end{alertblock}
\end{frame}

\begin{frame}
\frametitle{The Concept of Clustering}
\begin{center}
\includegraphics[width=\textwidth]{images/02_concept_clustering.png}
\end{center}
\end{frame}

\begin{frame}
\frametitle{What is Clustering?}
\begin{block}{Definition}
Clustering consists of dividing data into homogeneous groups where points within the same group resemble each other.
\end{block}

\vspace{1em}
\textbf{Goal:}
\begin{itemize}
    \item Points within the same cluster → similar
    \item Points from different clusters → different
\end{itemize}

\vspace{1em}
\textbf{No labels!}
\begin{itemize}
    \item The algorithm discovers the groups on its own
    \item The categories are not known in advance
    \item Exploration and pattern discovery
\end{itemize}
\end{frame}

\begin{frame}
\frametitle{Applications of Clustering}
\begin{center}
\includegraphics[width=\textwidth]{images/03_applications_clustering.png}
\end{center}
\end{frame}

\begin{frame}
\frametitle{Concrete Examples}
\begin{table}
\centering
\small
\begin{tabular}{lll}
\toprule
\textbf{Domain} & \textbf{Application} & \textbf{Goal} \\
\midrule
Marketing & Customer segmentation & Homogeneous groups \\
E-commerce & Recommendations & Similar customers \\
Biology & Gene groups & Similar functions \\
Documents & Organization & Similar topics \\
Security & Anomaly detection & Suspicious behavior \\
Images & Compression & Reducing colors \\
\bottomrule
\end{tabular}
\end{table}

\vspace{1em}
\begin{exampleblock}{Typical Use Case}
A company has 10,000 customers → create 4-5 segments for targeted marketing
\end{exampleblock}
\end{frame}

% ==========================================
% SECTION 2: K-MEANS
% ==========================================
\section{K-Means: The Main Algorithm}

\begin{frame}
\frametitle{K-Means: How Does It Work?}
\begin{center}
\includegraphics[width=\textwidth]{images/04_kmeans_principe.png}
\end{center}
\end{frame}

\begin{frame}
\frametitle{K-Means Algorithm - Steps}
\begin{block}{Principle}
K-Means finds K groups by iterating until convergence
\end{block}

\vspace{0.5em}
\textbf{Steps:}
\begin{enumerate}
    \item \textbf{Choose K}: Desired number of clusters
    \item \textbf{Initialize}: K random centers
    \item \textbf{Assignment}: Each point → nearest center
    \item \textbf{Update}: Recompute centers (mean of points)
    \item \textbf{Repeat}: Steps 3-4 until convergence
\end{enumerate}

\vspace{1em}
\begin{alertblock}{Convergence}
The centers stop moving = the algorithm is done!
\end{alertblock}
\end{frame}

\begin{frame}[fragile]
\frametitle{K-Means with Python}
\begin{lstlisting}[language=Python]
from sklearn.cluster import KMeans
import numpy as np

# Data
X = np.array([[1, 2], [1.5, 1.8], [5, 8], 
              [8, 8], [1, 0.6], [9, 11]])

# K-Means with K=2
kmeans = KMeans(n_clusters=2, random_state=42)

# Fit
kmeans.fit(X)

# Predictions (cluster labels)
labels = kmeans.labels_
print("Labels:", labels)  # [0, 0, 1, 1, 0, 1]

# Centers
centers = kmeans.cluster_centers_
print("Centers:", centers)
\end{lstlisting}
\end{frame}

\begin{frame}
\frametitle{Choosing K: The Elbow Method}
\begin{center}
\includegraphics[width=\textwidth]{images/05_methode_coude.png}
\end{center}
\end{frame}

\begin{frame}
\frametitle{How Do We Choose K?}
\begin{block}{Problem}
K-Means requires K to be specified in advance. How do we choose it?
\end{block}

\vspace{0.5em}
\textbf{Elbow Method:}
\begin{enumerate}
    \item Test several values of K (1, 2, 3, ..., 10)
    \item For each K, compute the inertia
    \item Plot Inertia vs K
    \item Look for the "elbow" (inflection point)
\end{enumerate}

\vspace{0.5em}
\textbf{Inertia:}
\begin{itemize}
    \item Sum of squared distances between points and their centers
    \item The lower the inertia, the better
    \item Always decreases as K increases
\end{itemize}

\vspace{0.5em}
\begin{alertblock}{The Elbow}
The point where inertia stops decreasing quickly = optimal K
\end{alertblock}
\end{frame}

\begin{frame}[fragile]
\frametitle{Elbow Method - Code}
\begin{lstlisting}[language=Python]
from sklearn.cluster import KMeans
import matplotlib.pyplot as plt

# Test K from 1 to 10
inertias = []
K_range = range(1, 11)

for k in K_range:
    kmeans = KMeans(n_clusters=k, random_state=42)
    kmeans.fit(X)
    inertias.append(kmeans.inertia_)

# Plot
plt.plot(K_range, inertias, 'bo-')
plt.xlabel('Number of clusters K')
plt.ylabel('Inertia')
plt.title('Elbow Method')
plt.show()

# Choose K at the elbow
# For example: K=3 or K=4
\end{lstlisting}
\end{frame}

\begin{frame}
\frametitle{K-Means: Advantages and Limits}
\begin{columns}
\begin{column}{0.48\textwidth}
\textbf{✓ Advantages}
\begin{itemize}
    \item Simple and fast
    \item Works well on large datasets
    \item Easy to understand
    \item Effective for spherical clusters
\end{itemize}
\end{column}

\begin{column}{0.48\textwidth}
\textbf{✗ Limits}
\begin{itemize}
    \item Must choose K in advance
    \item Sensitive to initialization
    \item Assumes spherical clusters
    \item Sensitive to outliers
\end{itemize}
\end{column}
\end{columns}

\vspace{1em}
\begin{block}{When to Use It?}
K-Means is excellent when we roughly know the number of groups and the clusters are round-shaped.
\end{block}
\end{frame}

% ==========================================
% SECTION 3: HIERARCHICAL CLUSTERING
% ==========================================
\section{Hierarchical Clustering}

\begin{frame}
\frametitle{Hierarchical Clustering}
\begin{center}
\includegraphics[width=\textwidth]{images/06_clustering_hierarchique.png}
\end{center}
\end{frame}

\begin{frame}
\frametitle{Principle of Hierarchical Clustering}
\begin{block}{Idea}
Build a hierarchy of clusters by progressively merging points/groups
\end{block}

\vspace{0.5em}
\textbf{Algorithm (bottom-up approach):}
\begin{enumerate}
    \item Initially: each point = 1 cluster
    \item Merge the 2 closest clusters
    \item Repeat until there is only 1 cluster
    \item Result: a tree (dendrogram)
\end{enumerate}

\vspace{0.5em}
\textbf{Main advantage:}
\begin{itemize}
    \item \textbf{No need to choose K in advance!}
    \item We "cut" the tree at the desired level
    \item Clear visualization with the dendrogram
\end{itemize}
\end{frame}

\begin{frame}
\frametitle{Dendrogram}
\textbf{What is a dendrogram?}
\begin{itemize}
    \item A tree showing how clusters merge
    \item Y-axis = distance between clusters
    \item Cutting the tree horizontally → choosing K
\end{itemize}

\vspace{1em}
\textbf{Reading it:}
\begin{itemize}
    \item Merge height = dissimilarity
    \item The higher the merge, the more different the groups
    \item Cut low → many clusters
    \item Cut high → few clusters
\end{itemize}
\end{frame}

\begin{frame}[fragile]
\frametitle{Hierarchical Clustering - Code}
\begin{lstlisting}[language=Python]
from sklearn.cluster import AgglomerativeClustering
from scipy.cluster.hierarchy import dendrogram, linkage
import matplotlib.pyplot as plt

# Data
X = [[1, 2], [1.5, 1.8], [5, 8], [8, 8]]

# Hierarchical clustering (K=2)
hier = AgglomerativeClustering(n_clusters=2)
labels = hier.fit_predict(X)
print("Labels:", labels)

# Dendrogram
linkage_matrix = linkage(X, method='ward')
plt.figure(figsize=(10, 5))
dendrogram(linkage_matrix)
plt.title('Dendrogram')
plt.xlabel('Points')
plt.ylabel('Distance')
plt.show()
\end{lstlisting}
\end{frame}

\begin{frame}
\frametitle{Hierarchical: Advantages and Limits}
\begin{columns}
\begin{column}{0.48\textwidth}
\textbf{✓ Advantages}
\begin{itemize}
    \item No need for K
    \item Informative dendrogram
    \item Natural hierarchy
    \item Deterministic
\end{itemize}
\end{column}

\begin{column}{0.48\textwidth}
\textbf{✗ Limits}
\begin{itemize}
    \item Slow on large datasets
    \item No reassignment
    \item Sensitive to noise
    \item O(n³) complexity
\end{itemize}
\end{column}
\end{columns}

\vspace{1em}
\begin{block}{When to Use It?}
Good for small/medium datasets when you want to explore the natural hierarchy of the data
\end{block}
\end{frame}

% ==========================================
% SECTION 4: DBSCAN
% ==========================================
\section{DBSCAN}

\begin{frame}
\frametitle{DBSCAN: Arbitrary Shapes}
\begin{center}
\includegraphics[width=\textwidth]{images/07_dbscan.png}
\end{center}
\end{frame}

\begin{frame}
\frametitle{DBSCAN: Principle}
\begin{block}{Definition}
DBSCAN (Density-Based Spatial Clustering) groups points in dense regions and identifies noise
\end{block}

\vspace{0.5em}
\textbf{Two parameters:}
\begin{itemize}
    \item \textbf{eps}: Neighborhood radius
    \item \textbf{min\_samples}: Minimum points to be considered dense
\end{itemize}

\vspace{0.5em}
\textbf{Types of points:}
\begin{itemize}
    \item \textbf{Core points}: At least min\_samples neighbors within eps
    \item \textbf{Border points}: Fewer than min\_samples neighbors, but near a core point
    \item \textbf{Noise}: Neither core nor border → outliers
\end{itemize}
\end{frame}

\begin{frame}
\frametitle{Advantages of DBSCAN}
\textbf{✓ Strengths:}
\begin{itemize}
    \item Detects arbitrary shapes (not just round ones!)
    \item Automatically identifies noise
    \item No need to specify K
    \item Robust to outliers
\end{itemize}

\vspace{1em}
\textbf{✗ Limitations:}
\begin{itemize}
    \item Sensitive to the eps and min\_samples parameters
    \item Struggles with varying densities
    \item Not efficient in high dimensions
\end{itemize}

\vspace{1em}
\begin{block}{When to Use It?}
Excellent for data with complex shapes and the presence of noise/outliers
\end{block}
\end{frame}

\begin{frame}[fragile]
\frametitle{DBSCAN - Code}
\begin{lstlisting}[language=Python]
from sklearn.cluster import DBSCAN
import numpy as np

# Data
X = np.array([[1, 2], [2, 2], [2, 3],
              [8, 7], [8, 8], [25, 80]])

# DBSCAN
dbscan = DBSCAN(eps=3, min_samples=2)
labels = dbscan.fit_predict(X)

print("Labels:", labels)
# Example: [0, 0, 0, 1, 1, -1]
# -1 = noise/outlier

# Number of clusters (excluding noise)
n_clusters = len(set(labels)) - (1 if -1 in labels else 0)
print(f"Number of clusters: {n_clusters}")
\end{lstlisting}
\end{frame}

% ==========================================
% SECTION 5: COMPARISON
% ==========================================
\section{Comparing the Algorithms}

\begin{frame}
\frametitle{Visual Comparison}
\begin{center}
\includegraphics[width=\textwidth]{images/08_comparaison_algorithmes.png}
\end{center}
\end{frame}

\begin{frame}
\frametitle{Comparison Table}
\begin{table}
\centering
\tiny
\begin{tabular}{lllll}
\toprule
\textbf{Algorithm} & \textbf{Needs K?} & \textbf{Shape} & \textbf{Speed} & \textbf{Use} \\
\midrule
K-Means & Yes & Spherical & Fast & Standard \\
Hierarchical & No & Any & Slow & Small dataset \\
DBSCAN & No & Arbitrary & Medium & Complex shapes \\
\bottomrule
\end{tabular}
\end{table}

\vspace{1em}
\textbf{Choose based on:}
\begin{itemize}
    \item \textbf{Data size}: K-Means (large), Hierarchical (small)
    \item \textbf{Cluster shape}: K-Means (round), DBSCAN (complex)
    \item \textbf{Knowledge of K}: K-Means (yes), others (no)
    \item \textbf{Presence of noise}: DBSCAN (yes), K-Means (no)
\end{itemize}
\end{frame}

% ==========================================
% SECTION 6: EVALUATION
% ==========================================
\section{Evaluating Clustering}

\begin{frame}
\frametitle{How Do We Evaluate a Clustering?}
\begin{alertblock}{Problem}
No labels → no accuracy! How do we know if it's good?
\end{alertblock}

\vspace{1em}
\textbf{Silhouette Score:}
\begin{itemize}
    \item Measures the quality of the clustering
    \item Ranges between -1 and +1
    \item For each point: compactness vs separation
\end{itemize}

\vspace{0.5em}
\textbf{Interpretation:}
\begin{itemize}
    \item \textbf{Close to +1}: Good clustering
    \item \textbf{Close to 0}: Points on the boundary
    \item \textbf{Negative}: Poor clustering
\end{itemize}
\end{frame}

\begin{frame}[fragile]
\frametitle{Silhouette Score - Code}
\begin{lstlisting}[language=Python]
from sklearn.metrics import silhouette_score
from sklearn.cluster import KMeans

# Data
X = [[1, 2], [1.5, 1.8], [5, 8], [8, 8], [1, 0.6], [9, 11]]

# Clustering
kmeans = KMeans(n_clusters=2, random_state=42)
labels = kmeans.fit_predict(X)

# Silhouette Score
score = silhouette_score(X, labels)
print(f"Silhouette Score: {score:.3f}")

# Interpretation
if score > 0.5:
    print("Good clustering")
elif score > 0.25:
    print("Acceptable clustering")
else:
    print("Poor clustering")
\end{lstlisting}
\end{frame}

\begin{frame}
\frametitle{Other Metrics}
\textbf{Internal metrics (without labels):}
\begin{itemize}
    \item \textbf{Silhouette}: Overall quality
    \item \textbf{Inertia}: Cluster compactness (K-Means)
    \item \textbf{Davies-Bouldin}: Compactness/separation ratio
\end{itemize}

\vspace{1em}
\textbf{External metrics (with labels, if available):}
\begin{itemize}
    \item \textbf{Adjusted Rand Index}: Similarity to true labels
    \item \textbf{Normalized Mutual Information}: Shared information
\end{itemize}

\vspace{1em}
\begin{block}{Advice}
Use Silhouette + visualization to evaluate
\end{block}
\end{frame}

% ==========================================
% SECTION 7: WORKFLOW
% ==========================================
\section{Practical Workflow}

\begin{frame}
\frametitle{Complete Workflow}
\begin{center}
\includegraphics[width=\textwidth]{images/09_workflow_complet.png}
\end{center}
\end{frame}

\begin{frame}
\frametitle{Clustering Pipeline}
\textbf{The 7 steps:}
\begin{enumerate}
    \item \textbf{Data}: Collection and exploration
    \item \textbf{Preprocessing}: Normalization (important!)
    \item \textbf{Choosing K}: Elbow method or other
    \item \textbf{Clustering}: Apply the algorithm
    \item \textbf{Evaluation}: Silhouette Score
    \item \textbf{Interpretation}: Analyze the groups
    \item \textbf{Action}: Use the insights
\end{enumerate}

\vspace{1em}
\begin{alertblock}{Normalization}
ALWAYS normalize the data before clustering!
\end{alertblock}
\end{frame}

\begin{frame}
\frametitle{Why Normalize?}
\begin{exampleblock}{Example Without Normalization}
\begin{itemize}
    \item Feature 1: Income (0 - €100,000)
    \item Feature 2: Age (0 - 100 years)
\end{itemize}
→ Income would dominate the distance!
\end{exampleblock}

\vspace{1em}
\textbf{Solution: StandardScaler}
\begin{itemize}
    \item Mean = 0
    \item Standard deviation = 1
    \item All features on the same scale
\end{itemize}

\vspace{1em}
\begin{block}{Rule}
Normalize BEFORE clustering, for K-Means and DBSCAN
\end{block}
\end{frame}

\begin{frame}[fragile]
\frametitle{Complete Template}
\begin{lstlisting}[language=Python]
from sklearn.preprocessing import StandardScaler
from sklearn.cluster import KMeans
from sklearn.metrics import silhouette_score
import pandas as pd

# 1. Load
df = pd.read_csv('data.csv')

# 2. Prepare
X = df[['feature1', 'feature2']].values

# 3. Normalize
scaler = StandardScaler()
X_scaled = scaler.fit_transform(X)

# 4. Choose K (elbow method)
# ... elbow code ...

# 5. Clustering
kmeans = KMeans(n_clusters=3, random_state=42)
labels = kmeans.fit_predict(X_scaled)

# 6. Evaluate
score = silhouette_score(X_scaled, labels)
print(f"Silhouette: {score:.3f}")

# 7. Interpret
df['Cluster'] = labels
print(df.groupby('Cluster').mean())
\end{lstlisting}
\end{frame}

% ==========================================
% SECTION 8: EXERCISES
% ==========================================
\section{Practice Exercises}

\begin{frame}
\frametitle{Exercise 1: Customer Segmentation}
\textbf{Dataset:} Customer data from an e-commerce site

\vspace{0.5em}
\textbf{Features:}
\begin{itemize}
    \item Annual income
    \item Spending score (1-100)
\end{itemize}

\vspace{1em}
\textbf{To do:}
\begin{enumerate}
    \item Load and visualize the data
    \item Normalize
    \item Use the elbow method to find K
    \item Apply K-Means
    \item Compute the Silhouette Score
    \item Interpret the segments (profiles)
\end{enumerate}
\end{frame}

\begin{frame}
\frametitle{Exercise 2: Comparing Algorithms}
\textbf{Dataset:} Scikit-learn's make\_moons (curved shapes)

\vspace{1em}
\textbf{To do:}
\begin{enumerate}
    \item Create the dataset with \texttt{make\_moons}
    \item Apply K-Means (K=2)
    \item Apply Hierarchical Clustering
    \item Apply DBSCAN
    \item Visualize the 3 results
    \item Compare visually: which works best?
    \item Compute Silhouette for each
\end{enumerate}

\vspace{0.5em}
\textbf{Goal:} Understand that DBSCAN is better for non-spherical shapes
\end{frame}

% ==========================================
% SECTION 9: BEST PRACTICES
% ==========================================
\section{Best Practices}

\begin{frame}
\frametitle{Clustering Checklist}
\begin{enumerate}
    \item ✓ \textbf{Always} normalize the data
    \item ✓ \textbf{Visualize} before clustering
    \item ✓ \textbf{Test} several values of K
    \item ✓ \textbf{Use} the elbow method
    \item ✓ \textbf{Compute} the Silhouette Score
    \item ✓ \textbf{Compare} several algorithms
    \item ✓ \textbf{Interpret} the clusters (business meaning)
    \item ✓ \textbf{Document} your choices
\end{enumerate}

\vspace{1em}
\begin{alertblock}{The Most Important Thing}
Interpretation! A good clustering must make sense for the domain.
\end{alertblock}
\end{frame}

\begin{frame}
\frametitle{Mistakes to Avoid}
\begin{block}{❌ DON'T}
\begin{itemize}
    \item Forget to normalize
    \item Choose K at random
    \item Only test one algorithm
    \item Ignore outliers
    \item Skip checking the Silhouette Score
    \item Cluster without visualizing
\end{itemize}
\end{block}

\vspace{1em}
\begin{exampleblock}{✅ DO}
\begin{itemize}
    \item Visually explore the data
    \item Normalize systematically
    \item Test K=2 to K=10 at minimum
    \item Compare K-Means, Hierarchical, DBSCAN
\end{itemize}
\end{exampleblock}
\end{frame}

% ==========================================
% SECTION 10: CONCLUSION
% ==========================================
\section{Conclusion}

\begin{frame}
\frametitle{What You Have Learned}
\textbf{Concepts:}
\begin{itemize}
    \item Difference between supervised/unsupervised
    \item Definition of clustering
    \item Concrete applications
\end{itemize}

\vspace{0.5em}
\textbf{Algorithms:}
\begin{itemize}
    \item \textbf{K-Means}: Fast, round clusters
    \item \textbf{Hierarchical}: Dendrogram, no need for K
    \item \textbf{DBSCAN}: Arbitrary shapes, detects noise
\end{itemize}

\vspace{0.5em}
\textbf{Practice:}
\begin{itemize}
    \item Elbow method
    \item Silhouette Score
    \item Complete workflow with Scikit-learn
\end{itemize}
\end{frame}

\begin{frame}
\frametitle{Key Points to Remember}
\begin{block}{1. Unsupervised = No Labels}
We discover hidden structures in the data
\end{block}

\vspace{0.5em}
\begin{block}{2. Normalization Is Mandatory}
Always use StandardScaler before clustering
\end{block}

\vspace{0.5em}
\begin{block}{3. Elbow Method}
For finding the optimal K in K-Means
\end{block}

\vspace{0.5em}
\begin{block}{4. Silhouette Score}
Measures quality: close to 1 = good
\end{block}

\vspace{0.5em}
\begin{block}{5. Interpretation}
The clusters must make business sense!
\end{block}
\end{frame}

\begin{frame}
\frametitle{When to Use Which Algorithm?}
\begin{table}
\centering
\small
\begin{tabular}{ll}
\toprule
\textbf{Situation} & \textbf{Algorithm} \\
\midrule
Large dataset, round clusters & K-Means \\
Small dataset, explore hierarchy & Hierarchical \\
Complex shapes, noise & DBSCAN \\
Don't know how many groups & Hierarchical or DBSCAN \\
Need for speed & K-Means \\
\bottomrule
\end{tabular}
\end{table}

\vspace{1em}
\begin{alertblock}{Advice}
In practice: test both K-Means AND DBSCAN, compare the results!
\end{alertblock}
\end{frame}

\begin{frame}
\frametitle{Resources}
\textbf{Documentation:}
\begin{itemize}
    \item Scikit-learn Clustering: \texttt{scikit-learn.org/clustering}
    \item Pandas, NumPy, Matplotlib
\end{itemize}

\vspace{1em}
\textbf{Datasets to practice:}
\begin{itemize}
    \item Iris (classic)
    \item Customer Segmentation (Kaggle)
    \item Make\_blobs, Make\_moons (Scikit-learn)
\end{itemize}

\vspace{1em}
\textbf{Next step:}
\begin{itemize}
    \item Model evaluation and improvement
    \item Advanced cross-validation
    \item Feature engineering
\end{itemize}
\end{frame}

\begin{frame}
\begin{center}
{\Huge Thank You!}

\vspace{2em}
{\Large Questions?}

\vspace{3em}
\textbf{Meyssem Bouzaien}\\
Academic Year 2025-2026

\vspace{2em}
\textit{Good luck with your clustering labs!}
\end{center}
\end{frame}

\end{document}
